% =============================================================================
% 02_ols_ridge_lasso.tex — Problem 2: OLS, Ridge, and Lasso
% =============================================================================
\section{OLS, Ridge, and Lasso Regression}
\label{sec:ols-ridge-lasso}

% ---- 2.1 OLS Baseline ------------------------------------------------------
\subsection{OLS Baseline (P2 — Tuấn)}
\label{sec:ols}

Ordinary Least Squares provides the unpenalised baseline:
\[
    \hat{\bbeta}^{\mathrm{OLS}}
    = \argmin_{\bbeta \in \R^p} \frac{1}{2n}\norm{\by - \bX \bbeta}_2^2.
\]

\input{../output/tables/tab_p2_ols_coef}

\includefigure{0.72\textwidth}{fig_p2_ols_resid}%
    {OLS residuals versus fitted values, with a LOESS trend.}%
    {fig:ols-residuals}

\includefigure{0.60\textwidth}{fig_p2_ols_qq}%
    {Normal Q-Q plot of the OLS residuals.}%
    {fig:ols-qq}

With standardised predictors, each slope is the change in body-fat percentage
points for a one-training-SD increase, holding the other measurements fixed.
\texttt{abdom} has the dominant positive slope (9.47), while \texttt{weight}
($-1.33$), \texttt{adipos} ($-1.32$), and \texttt{hip} ($-1.37$) have
counterintuitive negative partial slopes.  These signs should not be read
causally: they occur after conditioning on strongly correlated size measures.
OLS has training RMSE 3.884 and shared-fold CV RMSE 4.406.  The residual and
Q-Q plots show the fit diagnostics without using the holdout data; visible tail
departures also caution against relying only on Gaussian-error summaries.

% ---- 2.2 Ridge Regression --------------------------------------------------
\subsection{Ridge Regression (P2 — Tuấn)}
\label{sec:ridge}

Ridge regression adds an $\ell_2$ penalty to the OLS objective:
\[
    \hat{\bbeta}^{\mathrm{Ridge}}
    = \argmin_{\bbeta \in \R^p}
      \left\{
          \frac{1}{2n}\norm{\by - \bX \bbeta}_2^2
          + \frac{\lambda}{2} \norm{\bbeta}_2^2
      \right\}.
\]

% -- Ridge CV plot
\includefigure{0.75\textwidth}{fig_p2_ridge_cv}%
    {Ridge regression: cross-validated MSE as a function of $\log\lambda$.}%
    {fig:ridge-cv}

% -- Ridge coefficient path
\includefigure{0.75\textwidth}{fig_p2_ridge_path}%
    {Ridge coefficient path: coefficients vs.\ $\log\lambda$.}%
    {fig:ridge-path}

% -- Ridge coefficient bar chart
\includefigure{0.7\textwidth}{fig_p2_ridge_coef}%
    {Ridge coefficient estimates at $\lambda_{\min}$ vs.\ $\lambda_{1\mathrm{se}}$.}%
    {fig:ridge-coef-bar}

\paragraph{Tuning results.}
Shared-fold CV selected $\lambda_{\min}=0.6280$ (CV MSE 19.86, RMSE
4.457) and $\lambda_{1\mathrm{se}}=4.0367$ (CV MSE 23.38, RMSE 4.835).
The minimum rule targets the lowest estimated prediction error; the one-SE rule
accepts some estimated error for stronger shrinkage.

\input{../output/tables/tab_p2_ridge_lambda}
\input{../output/tables/tab_p2_ridge_coef}
\input{../output/tables/tab_p2_cond}

Ridge retains all 14 slopes but redistributes signal among correlated
measurements.  For example, \texttt{abdom} shrinks from 9.47 under OLS to 5.00
at $\lambda_{\min}$ and 2.15 at $\lambda_{1\mathrm{se}}$, while correlated
size measures receive smaller, more similar contributions.  On the
$G=\bX^\top\bX/n$ scale, adding $\lambda_{\min}I$ reduces the condition number
from 3166.75 to 15.17, a 208.7-fold improvement.  No Ridge slope is described
as an exact variable-selection result.

% =============================================================================
% SECTION 2.3: LASSO REGRESSION
% =============================================================================
\subsection{Lasso Regression (P3 — Thuận)}
\label{sec:lasso}

The Lasso replaces the $\ell_2$ penalty with an $\ell_1$ penalty, enabling
automatic feature selection:
\[
    \hat{\bbeta}^{\mathrm{Lasso}}
    = \argmin_{\bbeta \in \R^p}
      \left\{
          \norm{\by - \bX \bbeta}_2^2
          + \lambda \norm{\bbeta}_1
      \right\}.
\]

% -- Lasso CV plot
\includefigure{0.75\textwidth}{fig_p2_lasso_cv}%
    {Lasso regression: cross-validated MSE as a function of $\log\lambda$.}%
    {fig:lasso-cv}

% -- Lasso coefficient path
\includefigure{0.75\textwidth}{fig_p2_lasso_path}%
    {Lasso coefficient path: coefficients vs.\ $\log\lambda$.}%
    {fig:lasso-path}

\input{../output/tables/tab_p2_lasso_lambda}
\input{../output/tables/tab_p2_lasso_sel}
\input{../output/tables/tab_p2_lasso_coef}

At $\lambda_{\min}=0.0954$, the Lasso CV MSE is 17.99 (RMSE 4.241) and
nine slopes are nonzero: \texttt{age}, \texttt{weight}, \texttt{height},
\texttt{neck}, \texttt{abdom}, \texttt{hip}, \texttt{thigh},
\texttt{forearm}, and \texttt{wrist}.  At
$\lambda_{1\mathrm{se}}=0.7390$, only \texttt{age}, \texttt{height}, and
\texttt{abdom} remain and CV MSE rises to 20.39 (RMSE 4.516).  Unlike Ridge,
the $\ell_1$ optimality condition sets other slopes exactly to zero.  The
selection identifies predictive conditional associations, not causal effects.

% =============================================================================
% SECTION 2.4: TRAINING-BASED COMPARISON AND CORE MODEL
% =============================================================================
\subsection{Training-Based Comparison and Core Model (P3 — Thuận)}
\label{sec:comparison}

\input{../output/tables/tab_p2_comparison}

% -- Comparison bar chart
\includefigure{0.7\textwidth}{fig_p2_comparison}%
    {Cross-validated RMSE comparison: OLS vs.\ Ridge vs.\ Lasso.}%
    {fig:model-comparison}

\bigskip
\noindent
\fbox{\parbox{0.95\textwidth}{%
\textbf{Core Model Declaration (pre-holdout):}\\[4pt]
Based only on the shared five-fold training comparison, we nominate
\textbf{Lasso at $\lambda_{\min}=0.0954$}.  It has the lowest CV RMSE (4.241
versus 4.406 for OLS and 4.457 for Ridge) and reduces the model from 14 to nine
nonzero slopes.  This declaration precedes all holdout scoring.
}}

The $\lambda_{\min}$ rule answers which candidate minimises estimated
prediction error, whereas $\lambda_{1\mathrm{se}}$ asks for the most strongly
regularised model statistically close to that minimum.  Sparsity improves
deployment and communication but is not itself evidence of better prediction;
the CV error and active-set size therefore appear separately in the comparison.

% =============================================================================
% SECTION 2.5: PERTURBATION OR STABILITY CHECK
% =============================================================================
\subsection{Perturbation or Stability Check (P3 — Thuận)}
\label{sec:perturbation}

\paragraph{Prediction before check.}
Before resampling, we predicted that Ridge coefficients would vary less than
OLS coefficients, while Lasso would show selection instability among highly
correlated measurements.

% -- Perturbation plot
\includefigure{0.75\textwidth}{fig_p2_perturbation}%
    {Coefficient stability under predictor perturbation or bootstrap resampling.}%
    {fig:perturbation}

Using bootstrap seed \texttt{240302}, we drew $B=50$ training resamples and
refit the methods without consulting the holdout set.  The plot reports slope
standard deviations, and the Lasso table records selection frequencies.
Ridge is markedly more stable for the most collinear size variables: the
bootstrap SD for \texttt{weight} is 0.50 under Ridge versus 6.47 under OLS and
1.08 under Lasso.  Lasso consistently selects \texttt{abdom} and \texttt{wrist}
(frequency 1.00) and \texttt{forearm} (0.98), but alternates among
\texttt{weight} (0.46), \texttt{hip} (0.48), and \texttt{chest} (0.26), as
predicted for correlated predictors.

\input{../output/tables/tab_p2_lasso_stability}